\chapter{宠物基础数值模型}

本章的目标是回答一个基础但决定性的问題：玩家在界面中看到的宠物属性，究竟是如何由宠物的成长率生成的。若不先把这一层关系写清楚，后文关于转生、融合、蛋喂药与战斗表现的讨论就会失去共同基础。与很多玩家的直觉不同，宠物的 \vstd{} 并不会直接决定最终面板，而是先经过随机波动、档位判定、出生初始化、升级随机与显示公式，最后才表现为生命、攻击、防御与敏捷。

\section{宠物 \vstd{} 与成长率的关系}

在本文中，宠物的基础成长方向由四维向量
\begin{equation}
    \bm{g}=(V,S,T,D)
\end{equation}
表示，其中四个分量分别对应生命、攻击、防御与敏捷方向的成长率。对应的总和记为
\begin{equation}
    VSTD = V+S+T+D.
\end{equation}
需要强调的是，$VSTD$ 在研究中非常重要，但它本身并不是一个足以完全描述宠物的量。两只宠物即使具有相同的 $VSTD$，只要四维分配不同，其成长方向和实战表现也可能差异明显。因此，全文在大多数地方都同时保留“总和视角”和“分量视角”。

当前项目中，宠物原始成长点数到显示属性的换算可整理为
\begin{equation}
    HP = \left\lfloor \frac{4V_r + S_r + T_r + D_r}{100} \right\rfloor,
\end{equation}
\begin{equation}
    ATK = \left\lfloor \frac{S_r + 0.1T_r + 0.1V_r + 0.05D_r}{100} \right\rfloor,
\end{equation}
\begin{equation}
    DEF = \left\lfloor \frac{T_r + 0.1S_r + 0.1V_r + 0.05D_r}{100} \right\rfloor,
\end{equation}
\begin{equation}
    SPD = \left\lfloor \frac{D_r}{100} \right\rfloor,
    \label{eq:pet-display-map}
\end{equation}
其中 $(V_r,S_r,T_r,D_r)$ 代表进入显示公式之前的原始累计点数，而不是输入页面上的基础成长率本身。式~\eqref{eq:pet-display-map} 说明了四项成长率之间并非完全隔离：$V$ 主要影响生命，但也会少量进入攻击和防御；$D$ 主要影响敏捷，但同样会少量影响攻击和防御。因此，单纯以“主属性”理解宠物仍然不够，真正的映射关系是一个带交叉项的线性变换再叠加整数截断。

从数学上看，这一映射具有两个重要特点。第一，它是\emph{多对一}的：不同的原始点数组合，在通过向下取整之后，可能得到相同的显示属性。第二，它是\emph{层级式}的：显示属性不是直接由 $(V,S,T,D)$ 得到，而是由后续随机过程累积出的 $(V_r,S_r,T_r,D_r)$ 得到。也正因为如此，所谓“由显示属性反推成长率”天然只能得到候选集，而不能总是得到唯一解。

\section{等级、成长与初始化随机}

\subsection{一级初始化的三步结构}

对普通宠物而言，一级初始化并不是把基础成长率直接乘上某个系数，而是包含三步。第一步是四维独立的 $\pm2$ 波动：
\begin{equation}
    \varepsilon_i \in \{-2,-1,0,1,2\}, \qquad g_i^{(s)} = g_i + \varepsilon_i,
\end{equation}
其中 $g_i^{(s)}$ 表示真正存入成长档的“生效成长率”。第二步是用波动后的四项和判定 Rank，而不是用基础输入的 $VSTD$ 判定。第三步是在出生时进行一次总计 $10$ 点的随机分配，每次随机给四维中的某一项加 $1$。当前项目配套的宠物成长模拟，为了与既有模拟口径保持一致，在这一步额外采用了“单项最多额外得到 $4$ 点”的限制；这里需要特别说明，这一条是\emph{工具侧模拟设定}，不是本文已从服务器源码中直接验证出的硬编码上限。若按“$10$ 次独立、每次以 $1/4$ 概率落到某一维”的无上限模型看，对预先指定的一项有
\begin{equation}
    X \sim \mathrm{Binomial}(10, 1/4), \qquad
    \Pr(X \ge 5) = \sum_{k=5}^{10} \binom{10}{k} (1/4)^k (3/4)^{10-k} \approx 7.81\%.
\end{equation}
进一步地，$\Pr(X \ge 6) \approx 1.98\%$，$\Pr(X \ge 7) \approx 0.35\%$。因此，“单项超过 $4$ 点”在无上限随机模型下并非理论不可能，只是对某一指定维度而言概率不高。若将出生时随机加点记作 $e_i$，初始化参数记作 $INITNUM$，则一级原始点数满足
\begin{equation}
    g_{r,i}^{(1)} = INITNUM \cdot \bigl(g_i^{(s)} + e_i\bigr).
    \label{eq:lv1-raw}
\end{equation}
式~\eqref{eq:lv1-raw} 说明，一级面板已经同时混合了“基础成长率”“$\pm2$ 波动”“出生随机加点”和“初始化尺度”四类信息。

\subsection{0 转宠与转生宠的 Rank 判定差异}

Rank 是宠物成长系统中的关键中介变量，但它并不是对所有宠物共用同一张表。当前源码与工具口径下，至少要区分两类对象：其一是普通宠物，也就是玩家常说的“0 转宠”；其二是已经经过转生流程、重新进入成长路径的转生宠。两者虽然都会在后续成长中使用 Rank，但用于判定 Rank 的总和门槛并不相同。

对于普通宠物，Rank 由波动后的 $VSTD$ 判定，其分界如表~\ref{tab:normal-rank} 所示。

\begin{table}[H]
    \centering
    \caption{普通宠物（0 转宠）的 Rank 判定表}
    \label{tab:normal-rank}
    \begin{tabular}{cc}
        \toprule
        波动后 $VSTD$ & Rank \\
        \midrule
        $\ge 100$ & 0 \\
        $\ge 95$  & 1 \\
        $\ge 90$  & 2 \\
        $\ge 85$  & 3 \\
        $\ge 80$  & 4 \\
        $<80$      & 5 \\
        \bottomrule
    \end{tabular}
\end{table}

对于转生宠，Rank 判定门槛会发生变化，其分界如表~\ref{tab:trans-rank} 所示。

\begin{table}[H]
    \centering
    \caption{转生宠的 Rank 判定表}
    \label{tab:trans-rank}
    \begin{tabular}{cc}
        \toprule
        波动后 $VSTD$ & Rank \\
        \midrule
        $\ge 130$ & 0 \\
        $\ge 100$ & 1 \\
        $\ge 95$  & 2 \\
        $\ge 85$  & 3 \\
        $\ge 80$  & 4 \\
        $<80$      & 5 \\
        \bottomrule
    \end{tabular}
\end{table}

这两张表的差异意味着：同样是一个波动后的成长总和，放在 0 转宠和转生宠身上，不一定会落到相同 Rank。尤其是 Rank 0 的门槛从 $100$ 提高到 $130$，直接改变了高成长转生宠的判定标准。因此，任何把“普通宠的 Rank 规则”直接套用到转生宠上的分析，都会在后续成长倍率、随机区间和分布判断上产生系统性偏差。

从结构上说，Rank 的判定对象始终是“波动后的生效成长总和”，而不是输入页上的理想成长率总和；但判定所用的门槛表，会随着宠物是否已经转生而变化。后文第 4 章讨论转生模型时，会进一步说明：首转在计算原宠 Rank 时仍使用普通表，而转生后的新成长率在继续进入成长流程时，应按转生宠的 Rank 表理解。

Rank 一经确定，就会进入升级随机的倍率表。当前实现中，不同 Rank 对应的随机倍率区间如表~\ref{tab:frand-ranges} 所示。

\begin{table}[H]
    \centering
    \caption{不同 Rank 的升级倍率范围}
    \label{tab:frand-ranges}
    \begin{tabular}{cc}
        \toprule
        Rank & $fRand$ 对应整数区间 \\
        \midrule
        0 & $[450,500] \times 0.01$ \\
        1 & $[470,520] \times 0.01$ \\
        2 & $[490,540] \times 0.01$ \\
        3 & $[510,560] \times 0.01$ \\
        4 & $[530,580] \times 0.01$ \\
        5 & $[550,600] \times 0.01$ \\
        \bottomrule
    \end{tabular}
\end{table}

这组区间表明，Rank 并不是一个“纯评价标签”，而是后续升级过程中的随机参数入口。也就是说，同一只基础宠物如果因为 $\pm2$ 波动跨过了 Rank 边界，后续升级整条路径都会被改写。

\subsection{0 转宠与转生宠的完整成长口径对照}

玩家常说“0 转和转生后的宠物档位设置不同”。若按源码精确定义，这句话应理解为：两者真正不同的是 Rank 判定门槛及其进入点，而不是整个成长引擎都换了一套。一级初始化公式、出生 $10$ 点随机、升级时的 $10$ 点随机、$fRand$ 倍率表，以及最终显示公式，本质上仍是同一套结构；发生变化的是“用哪张 Rank 表解释波动后的成长总和”。这一点可整理为表~\ref{tab:normal-vs-trans-growth}。

\begin{table}[H]
    \centering
    \caption{0 转宠与转生宠的成长口径对照}
    \label{tab:normal-vs-trans-growth}
    \begin{tabularx}{\textwidth}{>{\raggedright\arraybackslash}p{3.6cm} >{\raggedright\arraybackslash}X >{\raggedright\arraybackslash}X}
        \toprule
        比较项目 & 0 转宠 & 转生宠 \\
        \midrule
        适用对象 & 普通捕捉宠、普通孵化宠、尚未进入转生路径的宠物 & 已完成转生并重新进入成长流程的宠物 \\
        Rank 判定对象 & 波动后的生效成长总和 $\sum_i g_i^{(s)}$ & 波动后的生效成长总和 $\sum_i g_i^{(s)}$ \\
        Rank 门槛 & 使用表~\ref{tab:normal-rank}：Rank 0 门槛为 $100$ & 使用表~\ref{tab:trans-rank}：Rank 0 门槛提高到 $130$ \\
        一级出生 $10$ 点随机 & 总计 $10$ 点；当前工具模拟默认额外采用单项最多 $4$ 点的限制 & 同左 \\
        升级时 $10$ 点随机 & 总计 $10$ 点，不设单项 $4$ 点限制 & 同左 \\
        $fRand$ 区间表 & 使用表~\ref{tab:frand-ranges} & 与 0 转宠相同，仍使用表~\ref{tab:frand-ranges} \\
        显示属性换算 & 使用式~\eqref{eq:pet-display-map} & 与 0 转宠相同，仍使用式~\eqref{eq:pet-display-map} \\
        \bottomrule
    \end{tabularx}
\end{table}

表~\ref{tab:normal-vs-trans-growth} 的意义在于，它把“档位设置不同”这件事拆成了“哪些变量真的改了，哪些步骤其实没改”。如果不做这一步区分，读者很容易误以为转生后连出生随机、升级随机甚至显示公式都一起改变；但从当前源码和工具实现看，这些部分并没有改动，真正变动的是 Rank 的解释规则，因此受影响最直接的是后续 $fRand$ 抽样区间的选择。

这里还需要额外区分两个容易混淆的 Rank 场景。第一，\emph{转生公式内部}为了计算 $F_x$，需要先给“原宠”判一次 Rank：首转使用普通表，二转使用转生表。第二，\emph{转生完成之后}，由转生公式生成的新成长率 $\bm{g}^{(new)}$ 再进入正常成长流程时，应把它视为“转生宠”，因此后续 $\pm2$ 波动后的 Rank 判定统一按转生表处理。前者服务于转生能力值公式本身，后者服务于转生后新宠的出生与升级随机，这两层 Rank 虽然名字相同，但作用阶段并不相同。

\subsection{升级随机模型}

在每次升级时，系统还会再做两件事。其一是重新进行一次总和为 $10$ 的随机加点，但这一次不再有“单项最多 $4$ 点”的限制；其二是根据 Rank 所对应的区间抽取一次倍率 $fRand$。若第 $\ell$ 级到第 $\ell+1$ 级的随机加点记为 $p_i^{(\ell)}$，则升级增量满足
\begin{equation}
    \Delta g_{r,i}^{(\ell)} = \left\lfloor \bigl(g_i^{(s)} + p_i^{(\ell)}\bigr) fRand^{(\ell)} \right\rfloor.
    \label{eq:levelup-delta}
\end{equation}
因此，高等级宠物的最终属性并不是某个固定点，而是一个分布。式~\eqref{eq:levelup-delta} 也解释了为什么项目中的成长模拟输出分位数而不是单值：在完整随机机制下，“最弱”“中位”“最强”本来就是不同样本路径的产物，而非同一路径中的不同观测角度。

\section{单项上限、总和上限与压缩机制}

从研究视角看，宠物系统中的上限可以分成两类。第一类是\emph{随机模型侧约束}，例如当前配套模拟在普通宠一级初始化时，对出生 $10$ 点随机分配额外施加的“单项最多 $4$ 点”限制；第二类是\emph{成长档压缩约束}，例如后文在蛋喂药、孵化和融合过程中会反复出现的单项上限与总和上限。前者主要影响初始点的离散分布，后者则直接改变成长档本身。

对普通成长模拟而言，最重要的结论是：出生随机加点不会改写已经存档的成长率，它只改变一级原始点数和之后的面板起点。与之相对，蛋喂药和孵化中的压缩会改变进入后续成长过程的成长档，因此它的影响会贯穿整个生命周期。也正因为此，后文在讨论养成优化时总是把“孵化后成长”视为真正的优化目标，而不是蛋基础成长或出生后一时的面板数值。

从数学上说，带上限的系统通常不是连续最优化问题，而是离散约束优化问题。无论是单项上限、总和上限，还是“先平均再乘系数”“先乘系数再取整”的顺序差异，都会造成可行域的折角与跳点。对宠物工具而言，这意味着搜索算法比闭式公式更重要；对论文而言，这意味着必须把上限发生的位置写清楚，否则所谓“某个方案更优”的结论就缺乏口径基础。

\section{蛋面数值与孵化后数值的关系}

融合与喂药系统引入了“蛋基础成长”和“孵化后成长”这两个容易混淆的概念。本文记蛋基础成长为
\begin{equation}
    \bm{b}=(b_V,b_S,b_T,b_D),
\end{equation}
它表示融合完成后、喂药前或喂药后的蛋面数值；再记孵化后成长为
\begin{equation}
    \bm{h}=(h_V,h_S,h_T,h_D),
\end{equation}
它表示按孵化公式折算之后、但尚未进入出生后 10 点随机加点之前的成长档。当前工具中，优化问题和大多数比较都以 $\bm{h}$ 为对象，而不是以一级显示属性为对象。

在无喂药情形下，当前实现中的确定性孵化模型可看作以 $\bm{b}$ 为输入、以 $\bm{h}$ 为输出的整数映射；而在有喂药情形下，还要先经过药效折算、单项裁切、总和压缩和按比例回分配等步骤，详见下一章。无论是否喂药，都必须区分“孵化后成长”和“出生后状态”：前者是进入成长流程的成长档，后者还会叠加后续的 $\pm2$ 波动、Rank 判定和出生随机加点。

因此，蛋面数值与孵化后数值之间的关系并不是“输入即输出”，而是“蛋面成长作为中间状态，经过一套带取整和上限的孵化映射后，生成真正参与后续成长的成长档”。从建模角度看，这一层正是把融合、喂药和普通成长统一起来的关键接口。后文关于转生、融合与喂药的所有公式，最终都会回到本章建立的两层结构：一层是成长档，另一层是由成长档驱动的随机成长过程。
